% Trabajo Final Individual — Ciencia de Datos 1C2026
% Formato pedido: Times New Roman 12, interlineado 1,5, márgenes normales, 6–8 páginas, APA.
% Motor: pdflatex
\documentclass[12pt]{article}

\usepackage[utf8]{inputenc}
\usepackage[T1]{fontenc}
\usepackage[spanish,es-noshorthands]{babel}
\usepackage{newtxtext,newtxmath} % Times New Roman
\usepackage[margin=1in]{geometry} % márgenes normales
\usepackage{setspace}
\onehalfspacing % interlineado 1,5
\usepackage{graphicx}
\usepackage{booktabs}
\usepackage{array}
\usepackage{amsmath}
\usepackage{microtype}
\usepackage{multicol}
\usepackage{xcolor}
\definecolor{linkblue}{HTML}{1A5FB4}
\usepackage[authoryear,round]{natbib}
\usepackage[colorlinks=true,linkcolor=black,citecolor=black,urlcolor=linkblue]{hyperref}
\usepackage{titlesec}
\titleformat{\section}{\normalfont\large\bfseries}{\thesection.}{0.5em}{}
\titleformat{\subsection}{\normalfont\normalsize\bfseries}{\thesubsection}{0.5em}{}
\titlespacing*{\section}{0pt}{0.5em}{0.2em}
\titlespacing*{\subsection}{0pt}{0.4em}{0.15em}
\setlength{\parskip}{0pt}
\setlength{\parindent}{1.4em}
\graphicspath{{experiments/analysis/figures/}{paper/figures/}}

\title{\vspace{-2.2em}\bfseries No todas las instrucciones se olvidan igual:\\
medición de la retención por tipo de instrucción\\ en asistentes de programación con LLMs}
\author{Tomás Pablo Korenblit\thanks{Una versión de este trabajo fue aceptada en el simposio ASAID, 55.\textsuperscript{as}~JAIIO (2026).}}
\date{Ciencia de Datos --- 1.\textsuperscript{er} cuatrimestre 2026}

\begin{document}
\maketitle
\vspace{-1.8em}

\begin{center}
\begin{minipage}{0.92\textwidth}
\small
\noindent\textbf{Resumen.} Los modelos de lenguaje grandes (LLMs) pierden
adherencia a las instrucciones a medida que la conversación se extiende. La
literatura mide ese fenómeno, pero siempre de forma agregada: informa una caída
promedio de cumplimiento agregando todas las instrucciones. Este trabajo aborda
una pregunta que esa literatura deja abierta: ¿decaen todas las instrucciones por
igual, o algunas se retienen y otras no? Para responder esto, construí un banco de
pruebas que planta doce preferencias de código como pedidos del usuario en
conversaciones sobre bases de código y mide si el modelo las respeta al final.
Sobre esas observaciones ajusté un modelo bayesiano con efectos jerárquicos por
tipo de decisión. El resultado principal es que la retención varía mucho entre
tipos de instrucción: los efectos del tratamiento cambian en más de un orden de
magnitud ($\sigma_\beta = 2{,}11$). Tres de doce instrucciones se benefician del
refuerzo (dos de ellas de forma robusta entre modelos), una empeora cuando se la
refuerza y ocho no requieren intervención, de modo que el refuerzo uniforme
malgasta la mayor parte del presupuesto. Esa heterogeneidad es lo que un sistema
de refuerzo selectivo necesitaría explotar.

\vspace{0.3em}
\noindent\textbf{Palabras clave:} cumplimiento de instrucciones; context rot;
inferencia bayesiana; refuerzo selectivo; evaluación de LLMs.
\end{minipage}
\end{center}

\vspace{0.5em}

\section{Introducción}

Cuando alguien usa un asistente conversacional como ChatGPT o Claude, la
experiencia inicial suele ser buena: explica su proyecto, pide que se respeten
ciertas reglas y el modelo obedece. Sin embargo, a medida que la conversación se
extiende, en algún momento el modelo deja de respetar lo que se le pidió: sigue
respondiendo, pero como si nunca hubiera recibido esas instrucciones. A este
fenómeno se lo denomina \emph{context rot}: la pérdida de adherencia a las
instrucciones conforme crece el contexto. El término abarca también fallas de
recuperación y razonamiento; acá lo uso en su acepción restringida a la adherencia
a instrucciones.

No se trata del defecto de un modelo particular; le pasa a todos, y las causas son
arquitectónicas. Una es la dilución de la atención: para generar cada token, un
LLM decide a qué partes del contexto asignarles más peso mediante un mecanismo de
atención que es un recurso finito que se reparte entre todos los tokens. Cuando el
contexto crece (las ventanas pasaron de decenas de miles de tokens en 2023 a un
millón desde principios de 2024, aproximadamente dos mil páginas), las
instrucciones del comienzo compiten por atención con todo lo posterior. Otra es el
sesgo posicional: la información del medio del contexto pesa menos que la de los
extremos \citep{liu2024}. La degradación, además, no aparece solo cerca del límite
de la ventana: se observa en todos los largos de contexto.

El problema importa por dos razones. Una es económica: las mitigaciones existentes
consisten en re-inyectar las instrucciones en cada turno, y eso multiplica el
costo en tokens a escala de miles de millones por día. El otro es la seguridad:
los \emph{guardrails} implementados a nivel de prompt (las reglas que impiden que
el modelo produzca contenido dañino) son instrucciones como cualquier otra, y se
sabe que los contextos largos aumentan sistemáticamente la tasa de jailbreak
\citep{anil2024}. Si decaen, el sistema no emite ningún error: falla en silencio.
Este trabajo aporta una estimación del efecto de tratamiento por tipo de
instrucción ---con un modelo jerárquico que cuantifica la heterogeneidad de la
retención--- y la justificación que esa medición ofrece para una estrategia de
refuerzo selectivo.

\section{Contexto y definición del problema}

Que los LLMs olvidan ya está documentado; lo que el trabajo busca entender es
\emph{cómo} lo hacen. La literatura informa métricas agregadas: un cumplimiento
medio que cae de un valor a otro. Pero un promedio oculta su estructura interna.
Si cada instrucción decae a un ritmo distinto, el promedio mezcla instrucciones
que el modelo retiene perfecto con otras que ignora.

Esa distinción determina qué estrategia de mitigación tiene sentido. Una
instrucción que el modelo ya cumple por defecto no necesita recordatorio; una que
empeora al repetirse no debería reforzarse jamás. Sin saber a qué tipo pertenece
cada una, toda política de refuerzo opera a ciegas. De aquí surge la pregunta que
guía el trabajo: \emph{¿la retención de instrucciones es homogénea o heterogénea
entre tipos de instrucción?} Si fuera homogénea, reforzar todo por igual bastaría;
si fuera heterogénea, hay que reforzar selectivamente, y para eso primero hay que
medirla.

\section{Revisión de la literatura}

Agrupo la literatura en cuatro bloques, según la función que cumplen en el
argumento.

\paragraph{Medición del olvido.} \citet{laban2025} compararon el mismo modelo con
toda la información de una vez frente a la misma tarea fragmentada en turnos, sobre
quince modelos, y hallaron una caída promedio del 39\% en el escenario multi-turno. Un detalle relevante: la \emph{aptitud} del modelo baja
poco, mientras que la inestabilidad sube mucho; lo que se degrada no es tanto la
capacidad como la fiabilidad. \citet{he2024}, en Meta, definieron el
\emph{Instruction Forgetting Ratio} y mostraron que la precisión de o1-preview cae
de 87,7\% a 70,7\% en tres turnos. \citet{du2025} agregaron que ampliar la ventana
de contexto no resuelve la degradación. Los tres informan promedios, pero ninguno
estima efectos por instrucción ni propone una política de refuerzo.

\paragraph{Mecanismo: posición y saturación.} \citet{liu2024} identificaron una
curva en forma de U (\emph{lost in the middle}): la información ubicada al comienzo
y al final pesa más que la del medio. Por eso las del system prompt (posición
cero) resisten mejor que las preferencias de mitad de conversación. \citet{mu2025}
describieron un efecto distinto, de saturación: al acumular reglas (de 1 a 20 en su
prueba), la adherencia a cada una cae marcadamente. Uno es de posición; el otro, de
cantidad, y juntos muestran que la adherencia no decae de modo uniforme sino
estructurado.

\paragraph{Implicaciones de seguridad.} \citet{anil2024} midieron empíricamente
que la tasa de \emph{jailbreak} aumenta de forma sistemática con la cantidad de
ejemplos previos en el contexto (\emph{many-shot jailbreaking}): cuanto más largo,
más vulnerable. \citet{rivasseau2025} aportó la formalización (el cociente entre
los tokens del system prompt y el contexto total tiende a cero a medida que la
conversación crece) y propuso interrupciones periódicas, aún sin evaluación
empírica.

\paragraph{Modelos bayesianos del comportamiento.} \citet{zhang2025} mostraron que
los LLM se comportan como filtros bayesianos descontados, con un factor
$\gamma<1$: el prior pierde peso frente a las observaciones recientes. Esto conecta
con \citet{laban2025}: una varianza que explota es justo lo que produce un
posterior inestable. En educación, el \emph{Bayesian Knowledge Tracing}
\citep{corbett1994} estima desde hace tres décadas, turno a turno, la probabilidad
de que un estudiante domine un concepto. La inversión conceptual que propongo es
tratar al LLM como el ``estudiante'' y a sus instrucciones como los ``conceptos''.

\section{Comparación de enfoques}

Las estrategias actuales comparten un mismo límite (Tabla~\ref{tab:mitig}).

\begin{table}[h]
  \centering
  \small
  \renewcommand{\arraystretch}{1.15}
  \caption{Estrategias de mitigación del \emph{context rot}.}
  \label{tab:mitig}
  \begin{tabular}{@{}>{\raggedright\arraybackslash}p{2.9cm} >{\raggedright\arraybackslash}p{4.2cm} >{\raggedright\arraybackslash}p{2.15cm} >{\raggedright\arraybackslash}p{3.05cm}@{}}
    \toprule
    \textbf{Enfoque} & \textbf{Idea} & \textbf{Costo} & \textbf{Límite} \\
    \midrule
    Repetir el prompt & Re-inyectar todas las reglas cada turno & $N\times$ entrada & Saturación: más reglas pueden empeorar \citep{mu2025} \\
    Jerarquía de instrucciones \citep{wallace2024} & Priorizar por origen, aprendido en el entrenamiento & Entrenamiento & Ordena roles, no instrucciones; es estática \\
    Duplicar el prompt \citep{leviathan2025} & Repetir la consulta (\texttt{<Q><Q>}) & $2\times$ entrada & Duplica todo por igual, sin priorizar \\
    Recordatorios fijos \citep{dongre2025} & Re-inyectar el objetivo en turnos predefinidos & $\sim$1,2$\times$ & Calendario fijo, no adaptativo \\
    \bottomrule
  \end{tabular}
\end{table}

Las cuatro hacen lo mismo: tratan al conjunto de instrucciones como un bloque y lo
refuerzan por igual. Ninguna decide \emph{qué} reforzar ni \emph{cuándo}. La
jerarquía de \citet{wallace2024} ordena roles, pero dentro de un mismo rol no
aborda el decaimiento temporal; la duplicación de \citet{leviathan2025} y los
recordatorios de \citet{dongre2025} reinyectan el contenido completo sin distinguir
entre lo que el modelo ya retuvo y lo que olvidó. Todas asumen, implícitamente, que
las instrucciones son intercambiables; la medición que sigue muestra que no lo son.

\section{Limitaciones de los enfoques actuales}

Algunos trabajos empiezan a distinguir tipos: \citet{longgenbench} muestran que
ciertas restricciones (de rango, periódicas) se degradan más que las de paso único
en generación larga. Pero ninguno estima el efecto por tipo de instrucción ni
deriva de ahí una política de refuerzo, y ese es el hueco. La métrica agregada
esconde la estructura que importa: un cumplimiento medio del 70\% no distingue
diez reglas al 100\% y diez al 40\% de veinte al 70\%, y esa diferencia cambia qué
conviene reforzar. El refuerzo uniforme, en consecuencia, desperdicia presupuesto
en instrucciones que no lo necesitan y puede ser contraproducente en aquellas que
empeoran con la repetición. Además, hasta donde llega esta revisión, las grandes
empresas y laboratorios (Anthropic, OpenAI, Google, Meta) no publicaron mediciones
de \emph{context rot} desagregadas por tipo de instrucción sino que las
mitigaciones son uniformes, en parte, porque no se cuenta con datos para
discriminar.

Esta limitación motiva directamente el trabajo: sin una medición por instrucción,
el refuerzo selectivo sería pura especulación. Por eso, antes de proponerlo, había
que producir la medición que lo justificara. La propuesta que sigue se deriva de
cubrir ese vacío.

\section{Propuesta de solución}

La propuesta tiene dos partes. La primera mide la heterogeneidad de la retención
(el grueso del trabajo y el aporte metodológico); la segunda es el sistema de
refuerzo selectivo que ese resultado habilita; el sistema solo se justifica si la
medición muestra esa heterogeneidad.

\subsection{El \emph{harness}: cómo medí la retención}

Para medir si una instrucción se olvida hace falta una conversación realista,
larga, con presión de contexto creciente, y al final una consigna que revele si el
modelo todavía respeta algo que se le pidió treinta mensajes atrás. Construí en
Python un \emph{harness} con las siguientes piezas.

\textbf{Codebases reales como fuente de presión.} En lugar de rellenar el contexto
con texto sintético, el harness opera sobre tres librerías reales del ecosistema
bayesiano de Python, elegidas por tamaño creciente: Bambi \citep{bambi2022}
($\sim$10K líneas), ArviZ \citep{arviz2019} ($\sim$25K) y PyMC \citep{pymc2023}
($\sim$57K). Durante los primeros veinte turnos inyecta en la conversación
archivos de código fuente reales de esas librerías, lo que genera un gradiente de
presión (Bambi alcanza $\sim$98K tokens de contexto; PyMC, $\sim$203K) y separa el
efecto del contexto del de la instrucción.

\textbf{Plantado de decisiones.} Definí doce ``decisiones'' de programación:
preferencias de estilo, arquitectura, testing, nomenclatura, dependencias y
documentación. En la condición \emph{tratamiento}, esas doce preferencias se
plantan como pedidos casuales, en línea, dentro de los mensajes del usuario (seis
temprano, en los turnos 1--7; seis tarde, en los 14--19); por ejemplo:
\emph{``Importante: cuando escribas operaciones con arrays en este proyecto, usá
siempre broadcasting de numpy en vez de for loops''}. Son preferencias del usuario
inyectadas a mitad de conversación, no reglas del system prompt: el caso difícil,
sin la ventaja posicional de la posición cero. La condición \emph{baseline} es
idéntica pero no planta nada, y mide qué hace el modelo por defecto.

\textbf{Turnos de evaluación.} En los turnos 20--25, el usuario solicita tareas de
generación de código diseñadas para forzar cada decisión. Cada turno de test
evalúa exactamente dos decisiones, una temprana y una tardía, para medir el olvido
reciente y también el de largo plazo.

\textbf{Verificadores determinísticos.} Aquí reside la reproducibilidad. Cada una
de las doce decisiones se evalúa con un \emph{checker} que analiza el código
generado mediante expresiones regulares y \emph{pattern matching}, y produce un
puntaje ordinal de 0 a 3 (0 = ignorada; 3 = cumplida por completo). No interviene
ningún LLM como juez en el \emph{scoring}: el resultado es reproducible entre
corridas. Tiene un límite de validez de constructo: cuando el checker no halla la
señal asigna 2 por defecto, y cuatro decisiones superan el 50\% de esos casos; las
tres con efectos más fuertes (\texttt{parametrize}, \texttt{standalone},
\texttt{no\_new}) tienen 0\% de ese caso, así que los hallazgos centrales no
dependen del artefacto. Por ejemplo, el verificador de \texttt{testing\_parametrize}
busca el decorador \texttt{@pytest.mark.parametrize} y penaliza funciones de test
separadas; el de \texttt{broadcasting} detecta bucles \texttt{for} frente a numpy.

\textbf{Lazo de herramientas y multiproveedor.} El \emph{runner} gestiona un ciclo
de llamadas a herramientas (leer, buscar, escribir con \texttt{write\_file}) sobre
una capa de proveedores intercambiable, que permitió evaluar cinco modelos vía
OpenRouter y la API de OpenAI. Como base, el harness usa DeepEval \citep{deepeval}
para las métricas con LLM como juez (persona y seguridad); en cambio, los doce
verificadores de decisión que sostienen estos resultados son propios y
determinísticos.

\subsection{Datos y modelo bayesiano}

El conjunto final consta de 28 conversaciones (9 \emph{baseline}, 19 tratamiento)
y 244 observaciones de cumplimiento (cada una es una decisión puntuada en un turno
de test) sobre cinco modelos (Qwen 3.5-27B como principal, con ambas condiciones;
más Gemini 3.1 Flash Lite, GPT-5.4-nano, Nemotron 3 Super 120B y Gemma 4, solo en
tratamiento).

Sobre esas observaciones ordinales ajusté una regresión logística ordinal con
efectos jerárquicos por tipo de decisión. El modelo ordinal es la elección
correcta porque la escala 0--3 no tiene intervalos iguales: el salto de
``ignorada'' (0) a ``mínima'' (1) no equivale al de ``casi'' a ``completa''. Para
cada observación, el cumplimiento latente
$\eta_i = \alpha_{d[i]} + \beta_{d[i]}\cdot\text{tratamiento}_i$ combina un
intercepto por decisión $\alpha_d$ (lo que el modelo hace sin que le digan) y un
efecto de tratamiento $\beta_d$ (cuánto mejora al pedírselo). El parámetro de
interés es $\sigma_\beta$, el desvío estándar grupal de los efectos: un
$\sigma_\beta$ grande indica que distintas instrucciones responden muy distinto a
ser enunciadas: la heterogeneidad que el refuerzo selectivo necesita. El ajuste se
hizo en PyMC \citep{pymc2023} con NUTS \citep{hoffman2014} vía nutpie
\citep{nutpie2024}, más verificaciones predictivas.

\subsection{Resultados}

El modelo convergió de forma limpia (cero transiciones divergentes,
$\hat{R}\leq 1{,}01$, tamaño de muestra efectivo mínimo superior a 1.400) y las
verificaciones predictivas posteriores reproducen la distribución observada de
puntajes. El desvío estándar grupal de los efectos de tratamiento resultó
$\sigma_\beta = 2{,}11$ (intervalo de credibilidad del 94\%: $[1{,}06,\ 3{,}28]$):
la heterogeneidad existe y es grande. Los efectos individuales, en escala de
\emph{log-odds}, van de $-2{,}36$ a $+5{,}44$ (Figura~\ref{fig:forest}). El
posterior separa las doce decisiones en tres grupos (Tabla~\ref{tab:res}).

\begin{figure}[h]
  \centering
  \includegraphics[width=0.40\textwidth]{treatment_effects_forest.png}
  \caption{Efectos de tratamiento por decisión ($\beta_d$) con intervalo de
  credibilidad del 94\%. Barras oscuras: el intervalo excluye el cero; rayada:
  efecto negativo; claras: no concluyentes.}
  \label{fig:forest}
\end{figure}

\emph{Instrucciones que se benefician del refuerzo.} Tres decisiones tienen un
efecto cuyo intervalo del 94\% queda por completo sobre cero.
\texttt{testing\_parametrize} es la más fuerte ($\beta = 5{,}44$): los modelos casi
nunca usan \texttt{parametrize} por su cuenta (media \emph{baseline}: 0,00), pero
retienen la preferencia cuando se la enuncian (media tratamiento: 2,79).
\texttt{module\_constants} ($\beta = 2{,}76$) y \texttt{architecture\_standalone}
($\beta = 2{,}38$) siguen el mismo patrón: línea de base cercana a cero, retención
apreciable al pedirlo.

\emph{Una instrucción que empeora con el refuerzo.} \texttt{dependencies\_no\_new}
($\beta = -2{,}36$; $P(\beta<0)=0{,}98$) es la única decisión donde enunciar la
regla \emph{reduce} el cumplimiento: la media \emph{baseline} es 2,25 (los modelos
ya evitan agregar imports), pero pedir ``no agregues imports nuevos'' la baja a
0,77. Es, además, la única instrucción formulada como negación, y eso hace pensar
que el problema está en el marco (la negación) y no en el contenido.

\emph{Instrucciones que no requieren intervención.} Las ocho restantes tienen
intervalos que cruzan el cero; varias ya parten de un cumplimiento alto
(\texttt{broadcasting} y \texttt{numpy\_style}: 3,00), y reforzarlas malgasta tokens.

Por último, el modelo no detectó varianza apreciable atribuible a la arquitectura
ni al tamaño del codebase ($\sigma_{\text{modelo}}\approx 0{,}4$;
$\sigma_{\text{codebase}}\approx 0{,}3$, con intervalos que cruzan el cero); en el
modelo completo, 267 transiciones divergentes llevaron a quitar esos efectos, de
modo que esta comparación es indicativa y no concluyente, más aún porque la
condición baseline se corrió en un solo modelo. Con esa salvedad, los datos son
consistentes con que lo que determina la retención es la instrucción misma. El
patrón se sostiene al restringir el análisis a Qwen ($r = 0{,}80$), con una
excepción: \texttt{architecture\_standalone} no replica en ese subconjunto
($\beta$ pasa de $2{,}38$ a $-0{,}71$), de modo que las decisiones claramente
beneficiosas se reducen de tres a dos.

\begin{table}[h]
  \centering
  \small
  \renewcommand{\arraystretch}{1.15}
  \caption{Prioridades de refuerzo derivadas del posterior. La ganancia es el cambio
  esperado en $P(\text{puntaje}\geq 2)$ al reforzar.}
  \label{tab:res}
  \begin{tabular}{@{}lrl@{}}
    \toprule
    \textbf{Decisión} & \textbf{Ganancia} & \textbf{Acción} \\
    \midrule
    \texttt{testing\_parametrize} & $+0{,}78$ & Reforzar \\
    \texttt{module\_constants}    & $+0{,}51$ & Reforzar \\
    \texttt{architecture\_standalone} & $+0{,}34$ & Reforzar \\
    \texttt{naming\_underscore}, \texttt{assert\_almost}, \texttt{regex\_comments} & $+0{,}14$ a $+0{,}24$ & Incierto \\
    6 decisiones restantes & $<+0{,}04$ & Omitir \\
    \texttt{dependencies\_no\_new} & $-0{,}47$ & Nunca reforzar \\
    \bottomrule
  \end{tabular}
\end{table}

\subsection{El sistema que esto habilita}

La medición es el prerrequisito de un sistema de refuerzo selectivo basado en
Bayesian Knowledge Tracing. En el harness ya está prototipado un \emph{tracker}
bayesiano que mantiene una creencia gaussiana $\mathcal{N}(\mu,\sigma^2)$ sobre el
cumplimiento latente de cada instrucción, la actualiza con cada observación
(aproximación de Laplace sobre un \emph{Graded Response Model}) y la hace decaer
entre turnos ($\gamma<1$). Sobre esa creencia, una política puede reforzar
\emph{selectivamente} solo las instrucciones con menor probabilidad de
cumplimiento, hasta un presupuesto fijo de tokens: \texttt{parametrize} recibe
refuerzo, \texttt{no\_new\_imports} queda excluida, y las que ya se cumplen no
gastan presupuesto. La Tabla~\ref{tab:res} es, en esencia, esa regla de asignación.

\section{Conclusiones}

Este trabajo midió la retención de instrucciones por tipo en asistentes de
programación basados en LLMs y encontró que es marcadamente heterogénea: los
efectos del tratamiento varían en más de un orden de magnitud. Algunas preferencias
se retienen apenas se las enuncia. La mayoría ya se cumplen por defecto, así que
reforzarlas solo desperdicia tokens. Y hay una que empeora cuando se la refuerza:
la única formulada como negación. No se detectó varianza apreciable atribuible al
modelo ni al nivel de presión de contexto, aunque el diseño limita esa comparación:
los datos apuntan a que lo que decide si una instrucción se retiene es la
instrucción misma.

La literatura mide el olvido en agregado y nunca por instrucción; ese es el dato
que faltaba y que acá medí. La implicancia práctica es directa: el refuerzo
uniforme malgasta la mayor parte de su presupuesto y puede ser contraproducente,
mientras que una política derivada de la medición indica cuáles conviene reforzar y
cuáles dejar sin intervención, incluida la que no hay que reforzar nunca. Las
probabilidades de retención estimadas son el insumo que un sistema de Bayesian
Knowledge Tracing necesitaría para funcionar.

Los comentarios de los revisores de ASAID orientan las siguientes líneas. La más
inmediata es ampliar y blindar la medición: correr la condición baseline en todos
los modelos (acá quedó atada a uno) y sumar réplicas, y complementar los
verificadores determinísticos con un panel de jueces LLM de otras familias, ciegos
a la condición, para separar el sesgo de medición del efecto real. Una segunda
línea, sugerida en la revisión, afila la negación: reformular como negación las
instrucciones que sí se benefician aislaría si pesa el marco o el contenido.
Quedan, además, la medición temporal densa para estimar $\gamma$, la evaluación en
lazo cerrado frente al refuerzo uniforme, y la validación sobre conversaciones
reales en lugar del banco simulado.

\begin{spacing}{1.0}
\scriptsize
\raggedright
\setlength{\bibsep}{0.12em}
\begin{multicols}{2}
\begin{thebibliography}{99}
\setlength{\itemsep}{0pt}
\setlength{\parsep}{0pt}

\bibitem[Abril-Pla et al.(2023)]{pymc2023} Abril-Pla, O., Andreani, V., Carroll, C., Dong, L., Fonnesbeck, C. J., Kochurov, M., \ldots{} Zinkov, R. (2023). PyMC: A modern, and comprehensive probabilistic programming framework in Python. \emph{PeerJ Computer Science, 9}, e1516. \url{https://doi.org/10.7717/peerj-cs.1516}

\bibitem[Anil et al.(2024)]{anil2024} Anil, C., Durmus, E., Sharma, M., Benton, J., Kundu, S., Batson, J., \ldots{} Duvenaud, D. (2024). Many-shot jailbreaking. \emph{Advances in Neural Information Processing Systems, 37}.

\bibitem[Capretto et al.(2022)]{bambi2022} Capretto, T., Piho, C., Kumar, R., Westfall, J., Yarkoni, T., \& Martin, O. A. (2022). Bambi: A simple interface for fitting Bayesian linear models in Python. \emph{Journal of Statistical Software, 103}(15), 1--29. \url{https://doi.org/10.18637/jss.v103.i15}

\bibitem[Corbett \& Anderson(1994)]{corbett1994} Corbett, A. T., \& Anderson, J. R. (1994). Knowledge tracing: Modeling the acquisition of procedural knowledge. \emph{User Modeling and User-Adapted Interaction, 4}(4), 253--278. \url{https://doi.org/10.1007/BF01099821}

\bibitem[Dongre et al.(2025)]{dongre2025} Dongre, V., Rossi, R. A., Lai, V. D., Yoon, D. S., Hakkani-Tür, D., \& Bui, T. (2025). \emph{Drift no more? Context equilibria in multi-turn LLM interactions} (\href{https://arxiv.org/abs/2510.07777}{arXiv:2510.07777}). arXiv.

\bibitem[Du et al.(2025)]{du2025} Du, Y., Tian, M., Ronanki, S., Rongali, S., Bodapati, S. B., Galstyan, A., \ldots{} Peng, H. (2025). Context length alone hurts LLM performance despite perfect retrieval. En \emph{Findings of the ACL: EMNLP 2025} (pp. 23281--23298).

\bibitem[He et al.(2024)]{he2024} He, Y., Jin, D., Wang, C., Bi, C., Mandyam, K., Zhang, H., \ldots{} Wang, S. (2024). \emph{Multi-IF: Benchmarking LLMs on multi-turn and multilingual instructions following} (\href{https://arxiv.org/abs/2410.15553}{arXiv:2410.15553}). arXiv.

\bibitem[Hoffman \& Gelman(2014)]{hoffman2014} Hoffman, M. D., \& Gelman, A. (2014). The No-U-Turn sampler: Adaptively setting path lengths in Hamiltonian Monte Carlo. \emph{Journal of Machine Learning Research, 15}, 1593--1623. \url{https://jmlr.org/papers/v15/hoffman14a.html}

\bibitem[Ip(2025)]{deepeval} Ip, J. (2025). \emph{DeepEval: The LLM evaluation framework} [Software]. Confident AI. \url{https://github.com/confident-ai/deepeval}

\bibitem[Kumar et al.(2019)]{arviz2019} Kumar, R., Carroll, C., Hartikainen, A., \& Martin, O. A. (2019). ArviZ: A unified library for exploratory analysis of Bayesian models in Python. \emph{Journal of Open Source Software, 4}(33), 1143. \url{https://doi.org/10.21105/joss.01143}

\bibitem[Laban et al.(2025)]{laban2025} Laban, P., Hayashi, H., Zhou, Y., \& Neville, J. (2025). \emph{LLMs get lost in multi-turn conversation} (\href{https://arxiv.org/abs/2505.06120}{arXiv:2505.06120}). arXiv.

\bibitem[Leviathan et al.(2025)]{leviathan2025} Leviathan, Y., Kalman, M., \& Matias, Y. (2025). \emph{Prompt repetition improves non-reasoning LLMs} (\href{https://arxiv.org/abs/2512.14982}{arXiv:2512.14982}). arXiv.

\bibitem[Liu et al.(2024)]{liu2024} Liu, N. F., Lin, K., Hewitt, J., Paranjape, A., Bevilacqua, M., Petroni, F., \& Liang, P. (2024). Lost in the middle: How language models use long contexts. \emph{Transactions of the Association for Computational Linguistics, 12}, 157--173. \url{https://doi.org/10.1162/tacl_a_00638}

\bibitem[Mu et al.(2025)]{mu2025} Mu, N., Lu, J., Lavery, M., \& Wagner, D. (2025). \emph{A closer look at system prompt robustness} (\href{https://arxiv.org/abs/2502.12197}{arXiv:2502.12197}). arXiv.

\bibitem[Rivasseau(2025)]{rivasseau2025} Rivasseau, T. (2025). \emph{LLM reinforcement in context} (\href{https://arxiv.org/abs/2511.12782}{arXiv:2511.12782}). arXiv.

\bibitem[Seyboldt(2024)]{nutpie2024} Seyboldt, A. (2024). \emph{nutpie: A fast sampler for Bayesian posteriors} [Software]. \url{https://github.com/pymc-devs/nutpie}

\bibitem[Wallace et al.(2024)]{wallace2024} Wallace, E., Xiao, K., Leike, R., Weng, L., Heidecke, J., \& Beutel, A. (2024). \emph{The instruction hierarchy: Training LLMs to prioritize privileged instructions} (\href{https://arxiv.org/abs/2404.13208}{arXiv:2404.13208}). arXiv.

\bibitem[Wu et al.(2025)]{longgenbench} Wu, Y., Hee, M. S., Hu, Z., \& Lee, R. K.-W. (2025). \emph{LongGenBench: Benchmarking long-form generation in long context LLMs} (\href{https://arxiv.org/abs/2409.02076}{arXiv:2409.02076}). ICLR 2025.

\bibitem[Zhang et al.(2025)]{zhang2025} Zhang, J., Yang, J., \& Wang, K. (2025). \emph{Large language models as discounted Bayesian filters} (\href{https://arxiv.org/abs/2512.18489}{arXiv:2512.18489}). arXiv.

\end{thebibliography}
\end{multicols}
\end{spacing}

\end{document}
